% ── Resum (Català) ─────────────────────────────────────────────────────────────
\selectlanguage{catalan}
\chapter*{Resum}
\addcontentsline{toc}{chapter}{Resum (Català)}

El \term{Reconeixement Òptic de Música} (OMR) és el procés d'obtenir representacions
simbòliques llegibles per màquina a partir d'imatges de partitures. Malgrat els avenços
en el camp, la transcripció automàtica de partitures reals de baixa qualitat continua
essent un repte obert. Aquesta memòria presenta un sistema OMR orientat a línies de
melodia monofòniques de jazz del recull \textit{The Real Book}.

El sistema es basa en una xarxa neuronal recurrent convolucional amb pèrdua CTC
(\term{CRNN-CTC}), entrenada sobre dades sintètiques generades a partir del conjunt
PrIMuS. Per reduir la diferència entre imatges sintètiques i escaneigs reals, s'aplica
una etapa d'augmentació que simula artefactes fotogràfics habituals. El model s'avalua
amb la \term{Taxa d'Error de Símbols} (SER), assolint un \textbf{1,17\%} al conjunt de
prova sintètic, i un ajust fi sobre pàgines reals del Real Book millora el rendiment
sobre imatges fora de domini.

El sistema complet integra detecció de pentagrames, reconeixement de melodia, extracció
de símbols d'acords i correcció gramatical, i accepta fitxers PDF o imatges com a
entrada. Aquest treball demostra la viabilitat dels models CRNN-CTC per al domini
específic de les línies de jazz monofòniques i estableix les bases per a extensions
futures cap a la transcripció polifònica.

% ── Resumen (Español) ──────────────────────────────────────────────────────────
\selectlanguage{spanish}
\chapter*{Resumen}
\addcontentsline{toc}{chapter}{Resumen (Español)}

El \term{Reconocimiento Óptico de Música} (OMR) es el proceso de obtener
representaciones simbólicas legibles por máquina a partir de imágenes de partituras.
A pesar de los avances en el campo, la transcripción automática de partituras reales
de baja calidad sigue siendo un desafío abierto. Esta memoria presenta un sistema OMR
orientado a líneas de melodía monofónicas de jazz de la colección \textit{The Real Book}.

El sistema se basa en una red neuronal recurrente convolucional con pérdida CTC
(\term{CRNN-CTC}), entrenada sobre datos sintéticos generados a partir del conjunto
PrIMuS. Para reducir la brecha entre imágenes sintéticas y escáneres reales, se aplica
una etapa de aumento que simula artefactos fotográficos habituales. El modelo se evalúa
con la \term{Tasa de Error de Símbolos} (SER), alcanzando un \textbf{1,17\%} en el
conjunto de prueba sintético, y un ajuste fino sobre páginas reales del Real Book mejora
el rendimiento sobre imágenes fuera de dominio.

El sistema completo integra detección de pentagramas, reconocimiento de melodía,
extracción de símbolos de acorde y corrección gramatical, y acepta archivos PDF o
imágenes como entrada. Este trabajo demuestra la viabilidad de los modelos CRNN-CTC
para el dominio específico de las líneas de jazz monofónicas y sienta las bases para
extensiones futuras hacia la transcripción polifónica.

% ── Abstract (English) ─────────────────────────────────────────────────────────
\selectlanguage{english}
\chapter*{Abstract}
\addcontentsline{toc}{chapter}{Abstract (English)}

\term{Optical Music Recognition} (OMR) is the task of automatically converting images
of musical scores into machine-readable symbolic representations. Despite advances in
the field, transcribing real-world scans of informal or low-quality sources remains an
open challenge. This thesis presents an end-to-end OMR system targeting monophonic
melody lines in jazz lead sheets, with a focus on the \textit{Real Book} collection.

The system is built around a \term{Convolutional Recurrent Neural Network with
Connectionist Temporal Classification} loss (CRNN-CTC), trained on synthetic data
derived from the PrIMuS dataset. To bridge the gap between clean synthetic renderings
and real scans, an augmentation stage simulates common photographic artefacts.
Evaluated using the \term{Symbol Error Rate} (SER), the model achieves \textbf{1.17\%}
on the synthetic test set, and a fine-tuning stage on real Real Book pages further
improves performance on out-of-domain images.

The full inference pipeline integrates staff detection, melody recognition, chord symbol
extraction, and a grammar correction stage, accepting PDF or image input. This work
demonstrates the feasibility of CRNN-CTC-based OMR for the specific domain of
monophonic jazz lead sheets and lays the groundwork for future extensions towards
polyphonic transcription.
